\documentclass[11pt]{article}

% ============================================================
% FOUNDATIONAL CANONICAL SPECIFICATION
% MACHINE-VERIFIABLE LAW
% R3 / MDM FIXED-POINT REVISION
% ============================================================

\usepackage[margin=1in]{geometry}
\usepackage{amsmath,amssymb,amsthm,mathtools}
\usepackage{booktabs}
\usepackage{array}
\usepackage{longtable}
\usepackage{enumitem}
\usepackage{hyperref}
\usepackage[expansion=false]{microtype}
\usepackage[T1]{fontenc}
\IfFileExists{lmodern.sty}{\usepackage{lmodern}}{}

\hypersetup{
    colorlinks=true,
    linkcolor=black,
    citecolor=black,
    urlcolor=blue,
    pdftitle={Foundational Canonical Specification for Machine-Verifiable Law},
    pdfauthor={Universal Standard Axiom Corporation},
    pdfsubject={
        Machine-Verifiable Law, Recursive Pursuit of Absolute Truth,
        Meta-Dynamical Mathematics, Canonical Verification
    }
}

% ============================================================
% CANONICAL MACROS
% ============================================================

\newcommand{\Rthree}{\ensuremath{\mathsf{R}^{3}}}
\newcommand{\MDM}{\ensuremath{\mathrm{MDM}}}
\newcommand{\MVL}{\ensuremath{\mathrm{MVL}}}
\newcommand{\WAD}{\ensuremath{\mathrm{WAD\mbox{-}18}}}
\newcommand{\CSL}{\ensuremath{\mathrm{CSL}}}
\newcommand{\SAC}{\ensuremath{\mathrm{SAC}}}

\newcommand{\Purpose}{\mathcal{U}}
\newcommand{\Authority}{\mathcal{A}}
\newcommand{\Evidence}{\mathcal{E}}
\newcommand{\Inv}{\mathcal{I}}
\newcommand{\State}{\mathcal{S}}
\newcommand{\Law}{\mathcal{L}}
\newcommand{\Prov}{\mathcal{P}}
\newcommand{\Truth}{\mathcal{T}}
\newcommand{\Rule}{\mathcal{R}}
\newcommand{\Spec}{\mathcal{C}}
\newcommand{\Derive}{\mathsf{Derive}}
\newcommand{\Verify}{\mathsf{Verify}}
\newcommand{\Refine}{\mathsf{Refine}}
\newcommand{\Sem}{\mathsf{Sem}}
\newcommand{\Accept}{\mathsf{ACCEPT}}
\newcommand{\Reject}{\mathsf{REJECT}}
\newcommand{\Unknown}{\mathsf{UNVERIFIED}}
\newcommand{\Obligation}{\mathsf{OBLIGATION}}
\newcommand{\Fixed}{\mathsf{Fix}}
\newcommand{\Domain}{\mathcal{D}}
\newcommand{\CCIO}{\ensuremath{\mathrm{CCIO}}}

% ============================================================
% THEOREM ENVIRONMENTS
% ============================================================

\newtheorem{axiom}{Axiom}
\newtheorem{definition}{Definition}
\newtheorem{principle}{Principle}
\newtheorem{proposition}{Proposition}
\newtheorem{requirement}{Requirement}

% ============================================================
% TITLE
% ============================================================

\title{
    \textbf{Foundational Canonical Specification}\\
    \large Machine-Verifiable Law\\[0.5em]
    \normalsize
    Recursive Pursuit of Absolute Truth,
    Meta-Dynamical Mathematics,
    Invariants, Verification, and Canonical State Derivation
}

\author{
    Universal Standard Axiom Corporation\\
    R\&D Division\\
    Foundational Canonical Specification
}

\date{October 2026}

% ============================================================
% DOCUMENT
% ============================================================

\begin{document}

\maketitle

\begin{center}
\Large
\textbf{BEGIN WITH PURPOSE. END IN TRUTH.}
\end{center}

\vspace{1em}

\begin{center}
\fbox{
\parbox{0.90\textwidth}{
\textbf{Canonical Status.}
This document defines the canonical semantic specification for
Machine-Verifiable Law (\MVL{}). It establishes the definitions,
relationships, invariants, state model, verification semantics,
provenance requirements, refinement procedure, and derivation
rules from which authorized formal, computational, conversational,
executable, and cryptographic representations may be produced.

No derivative representation shall silently alter the semantic
content, authority structure, invariant structure, or verification
semantics established here.
}}
\end{center}

\tableofcontents
\newpage

% ============================================================
\section{Purpose and Governing Direction}
% ============================================================

The purpose of Machine-Verifiable Law (\MVL{}) is to establish a
general formal architecture through which authorized normative
intent, applicable rules, invariants, evidence, provenance, state,
and state transitions may be represented and evaluated by
machines in an explicit, reproducible, auditable, and
appropriately deterministic manner.

The governing direction is:

\[
\boxed{
\Purpose
\longrightarrow
\Rthree
\longrightarrow
\MDM
\longrightarrow
\Inv
\longrightarrow
\Law
\longrightarrow
\Verify
\longrightarrow
\Truth
}
\]

The architecture begins with purpose because computation does not
by itself establish what a result is intended to mean.

The architecture ends in truth because an internally consistent
computation is not sufficient to establish that the resulting
state satisfies the intended and authorized conditions.

Accordingly:

\[
\boxed{
\textbf{BEGIN WITH PURPOSE. END IN TRUTH.}
}
\]

This statement is the governing direction of the specification.

It is not merely descriptive language.

It establishes the order in which purpose, refinement,
mathematical representation, invariant formation, formalization,
verification, and truth-constrained state are related.

% ============================================================
\section{Scope}
% ============================================================

This specification defines:

\begin{enumerate}[label=\arabic*.]
    \item the semantic foundations of Machine-Verifiable Law;
    \item the canonical relationship between purpose, authority,
          evidence, \Rthree{}, \MDM{}, invariants, law, state,
          verification, and truth;
    \item \Rthree{} as Recursive Pursuit of Absolute Truth;
    \item \MDM{} as the structural mathematical reference layer;
    \item the representation of applicable law as formal constraints
          upon admissible states and transitions;
    \item the canonical state and verification models;
    \item provenance and reproducibility requirements;
    \item fail-closed semantics for unresolved mandatory conditions;
    \item cross-domain structural correspondence;
    \item canonical fixed-point and convergence semantics;
    \item semantic-preservation requirements for derivative artifacts;
    \item controlled derivation into formal, conversational,
          executable, and provenance representations.
\end{enumerate}

The specification is domain-independent.

A domain may provide its own substantive rules, variables,
authorities, evidence, and invariants while conforming to the
common structural architecture defined herein.

The common architecture does not imply identical substantive
rules across domains.

% ============================================================
\section{Canonical Semantic Root}
% ============================================================

\begin{definition}[Canonical Semantic Root]
The \emph{Canonical Semantic Root} is the authoritative semantic
specification from which all authorized derivative
representations are derived.
\end{definition}

The canonical specification is represented by:

\[
\boxed{
\Spec =
(
\Purpose,
\Authority,
\Evidence,
\Inv,
\Law,
\State,
\Verify,
\Prov,
v
)
}
\]

where:

\begin{align*}
\Purpose &= \text{declared purpose},\\
\Authority &= \text{authorized normative source and authority structure},\\
\Evidence &= \text{evidentiary basis},\\
\Inv &= \text{canonical invariant set},\\
\Law &= \text{formalized applicable law},\\
\State &= \text{defined state space},\\
\Verify &= \text{verification semantics},\\
\Prov &= \text{provenance structure},\\
v &= \text{canonical version}.
\end{align*}

The canonical specification is therefore not merely a document,
file, program, database record, or hash.

It is the semantic object represented by the controlled
canonical artifact.

A particular file format may encode that object, but the encoding
shall not be confused with the semantic root itself.

% ============================================================
\section{Canonical Identity}
% ============================================================

Every canonical object shall possess a stable identity.

Let:

\[
id(C)
\]

denote the canonical identity of specification object $C$.

A versioned canonical object is:

\[
C_v=(id,v,\Sem(C),\Prov(C)).
\]

The canonical identity shall remain stable across derivative
representations unless the canonical semantic object itself is
formally superseded.

A change to canonical meaning constitutes a canonical revision.

A change to representation without a change in canonical meaning
constitutes a derivative revision.

Therefore:

\[
\boxed{
\text{Representation Change}
\neq
\text{Semantic Change}
}
\]

unless semantic comparison establishes otherwise.

% ============================================================
\section{Foundational Definitions}
% ============================================================

\begin{definition}[Purpose]
Purpose is the explicitly declared intended objective that
establishes why a law, invariant, system, action, state, or
transition exists.
\end{definition}

\begin{definition}[Authority]
Authority is the legitimate source or authorization structure
under which a normative rule is established, adopted, modified,
or applied.
\end{definition}

\begin{definition}[Machine-Verifiable Law]
Machine-Verifiable Law is law represented with sufficient formal
structure that an authorized verification mechanism can evaluate
whether a specified state or transition satisfies the applicable
formal constraints within a declared verification domain.
\end{definition}

\begin{definition}[State]
A state is a formally representable configuration of relevant
entities, values, relationships, permissions, conditions,
evidence, and provenance at a specified point or transition.
\end{definition}

\begin{definition}[Invariant]
An invariant is a formally represented condition that must remain
satisfied for a state or transition to be admissible under the
applicable canonical specification.
\end{definition}

\begin{definition}[Verification]
Verification is the evaluation of a specified state or transition
against the formally defined conditions that determine its
verification status.
\end{definition}

\begin{definition}[R\textsuperscript{3}]
\Rthree{} denotes \emph{Recursive Pursuit of Absolute Truth}.

It is the recursive refinement operator by which purpose,
definitions, assumptions, evidence, provenance, rules,
invariants, states, representations, contradictions, and
conclusions are examined and reconciled until either a
truth-constrained fixed point is reached or one or more
mandatory conditions remain unresolved.
\end{definition}

\begin{definition}[Meta-Dynamical Mathematics]
\MDM{} is the mathematical reference framework used to represent
relationships among states, transformations, invariants,
observations, mappings, and fixed points across domains.
\end{definition}

\begin{definition}[Fixed Point]
For an operator $F$ acting upon a specified state space, a fixed
point is a state $X^{*}$ satisfying:

\[
F(X^{*})=X^{*}.
\]

For the canonical specification, an \Rthree{} fixed point is a
canonical state whose materially relevant semantic, logical,
structural, and verification properties remain unchanged under
the declared refinement operator.
\end{definition}

\begin{definition}[Provenance]
Provenance is the traceable record establishing the origin,
authority, version, transformation history, and evidentiary basis
of a canonical object, rule, state, or result.
\end{definition}

\begin{definition}[Reproducibility]
Reproducibility is the property that equivalent canonical inputs,
rules, versions, evidence, and declared execution conditions
produce the same specified verification result.
\end{definition}

\begin{definition}[Semantic Preservation]
Semantic preservation is the property that a derivative
representation retains the canonical meaning applicable to its
declared representational domain.
\end{definition}

% ============================================================
\section{Foundational Axioms}
% ============================================================

\begin{axiom}[Purpose Precedes Consequential Computation]
A consequential machine-verifiable process shall have an
explicitly declared purpose before its consequential computation
is accepted as meaningful.

\[
\boxed{
\Purpose \prec \text{Consequential Computation}
}
\]
\end{axiom}

\begin{axiom}[Recursive Pursuit of Absolute Truth]
The canonical refinement mechanism shall recursively examine
purpose, definitions, assumptions, evidence, provenance, rules,
invariants, state, representation, contradictions, and
conclusions.

\[
\boxed{
\Rthree{}=\text{Recursive Pursuit of Absolute Truth}
}
\]
\end{axiom}

\begin{axiom}[Explicit Invariants]
Conditions upon which admissibility depends shall be explicitly
represented rather than left solely to undocumented
implementation interpretation.
\end{axiom}

\begin{axiom}[Authority]
Normative meaning shall be traceable to an identified and
authorized source.
\end{axiom}

\begin{axiom}[Provenance]
Consequential claims, rules, states, transitions, and verification
results shall have traceable provenance sufficient for their
declared purpose.
\end{axiom}

\begin{axiom}[Reproducibility]
A canonical verification result shall be reproducible from the
canonical inputs and declared execution conditions applicable to
that verification.
\end{axiom}

\begin{axiom}[No Hidden Consequential State]
A consequential verification result shall not depend upon
undocumented hidden state.
\end{axiom}

\begin{axiom}[Fail Closed]
A mandatory condition that cannot be established shall not be
silently converted into acceptance.

\[
\boxed{
\Unknown\neq\Accept
}
\]
\end{axiom}

\begin{axiom}[Deterministic Verification]
Where deterministic verification is specified, equivalent
canonical inputs and execution conditions shall produce the same
verification result.
\end{axiom}

\begin{axiom}[Semantic Preservation]
Every conforming derivative representation shall preserve the
canonical meaning applicable to its declared scope.
\end{axiom}

\begin{axiom}[No Silent Substitution]
An implementation shall not silently substitute its own purpose,
rule, invariant, interpretation, authority, or verification
semantics for those established by the canonical specification.
\end{axiom}

\begin{axiom}[Auditability]
A consequential result shall be traceable through its state,
applicable rule, invariant, evidence, provenance, version, and
authorized source.
\end{axiom}

\begin{axiom}[Truth-Constrained Acceptance]
Acceptance requires satisfaction of every mandatory condition
within the declared verification domain.
\end{axiom}

% ============================================================
\section{Canonical Dependency Structure}
% ============================================================

The canonical dependency structure is:

\[
\boxed{
\Purpose
\rightarrow
\Authority
\rightarrow
\Evidence
\rightarrow
\Rthree
\rightarrow
\MDM
\rightarrow
\Inv
\rightarrow
\Law
\rightarrow
\State
\rightarrow
\Verify
\rightarrow
\Truth
}
\]

This ordering does not imply that every implementation executes
each component sequentially.

It establishes semantic dependency.

In particular:

\[
\text{Verification}
\]

cannot independently determine:

\[
\text{Purpose}
\]

or:

\[
\text{Normative Authority}.
\]

Those must be supplied by the canonical specification.

Likewise:

\[
\text{Deterministic Execution}
\]

cannot independently establish:

\[
\text{Semantic Correctness}.
\]

% ============================================================
\section{Formal Model of Machine-Verifiable Law}
% ============================================================

A canonical law object is represented as:

\[
\boxed{
\Law =
(
\Purpose,
\Authority,
\Evidence,
\Inv,
\State,
\delta,
\Verify,
\Prov,
v
)
}
\]

where:

\begin{align*}
\Purpose &= \text{declared purpose},\\
\Authority &= \text{normative authority},\\
\Evidence &= \text{evidentiary basis},\\
\Inv &= \text{mandatory and conditional invariants},\\
\State &= \text{state space},\\
\delta &= \text{authorized state-transition relation},\\
\Verify &= \text{verification function},\\
\Prov &= \text{provenance},\\
v &= \text{version}.
\end{align*}

The verification function is:

\[
\boxed{
\Verify:
\State\times\Law
\rightarrow
\{
\Accept,\Reject,\Unknown
\}.
}
\]

For a state $S$:

\[
\Verify(S,\Law)=\Accept
\]

means that the state satisfies all mandatory verification
conditions applicable within the declared domain.

Likewise:

\[
\Verify(S,\Law)=\Reject
\]

means that at least one mandatory condition is affirmatively
violated.

Finally:

\[
\Verify(S,\Law)=\Unknown
\]

means that one or more mandatory conditions cannot presently be
established as satisfied or violated.

The three states are semantically distinct.

% ============================================================
\section{State-Transition Semantics}
% ============================================================

Let:

\[
S_t\in\State
\]

be a current state and:

\[
S_{t+1}\in\State
\]

a proposed subsequent state.

The transition is represented by:

\[
S_t\xrightarrow{\delta}S_{t+1}.
\]

The transition is admissible only when:

\[
\boxed{
\Verify(S_t\rightarrow S_{t+1},\Law)=\Accept.
}
\]

The notation:

\[
\Verify(S_t\rightarrow S_{t+1},\Law)
\]

denotes verification of the transition itself, including all
state-dependent and transition-dependent conditions.

If a mandatory condition is affirmatively violated:

\[
\boxed{
\Verify(S_t\rightarrow S_{t+1},\Law)=\Reject.
}
\]

If a mandatory condition cannot be established:

\[
\boxed{
\Verify(S_t\rightarrow S_{t+1},\Law)=\Unknown.
}
\]

For a fail-closed consequential implementation:

\[
\boxed{
\Verify\neq\Accept
\Rightarrow
\text{consequential transition not authorized by MVL}.
}
\]

Any exception mechanism shall itself be part of the applicable
canonical law and shall therefore be subject to verification.

% ============================================================
\section{Recursive Pursuit of Absolute Truth}
% ============================================================

Let:

\[
C_0
\]

be an initial canonical specification state.

The recursive process is:

\[
C_0
\xrightarrow{\Rthree}
C_1
\xrightarrow{\Rthree}
C_2
\xrightarrow{\Rthree}
\cdots
\]

with:

\[
\boxed{
C_{n+1}=\Rthree(C_n).
}
\]

Each iteration shall examine, where applicable:

\begin{enumerate}
    \item purpose;
    \item authority;
    \item definitions;
    \item assumptions;
    \item evidence;
    \item provenance;
    \item logical consistency;
    \item invariant completeness;
    \item state representation;
    \item transition semantics;
    \item computational representation;
    \item reproducibility;
    \item semantic preservation;
    \item contradictions;
    \item unresolved obligations.
\end{enumerate}

The recursive process shall produce a revised state whenever a
material defect is identified.

A material defect includes, at minimum:

\begin{enumerate}
    \item contradiction;
    \item semantic ambiguity affecting interpretation;
    \item unsupported assertion presented as established;
    \item missing mandatory dependency;
    \item inconsistent mathematical typing;
    \item inconsistent verification semantics;
    \item hidden consequential state;
    \item semantic drift between representations;
    \item unresolved mandatory obligation incorrectly treated as
          satisfied.
\end{enumerate}

% ============================================================
\section{R3 Refinement Operator}
% ============================================================

The \Rthree{} operator is defined structurally as:

\[
\boxed{
\Rthree:
\mathfrak{C}\rightarrow\mathfrak{C}
}
\]

where $\mathfrak{C}$ is the space of canonical specification
states.

For a canonical state $C$:

\[
\Rthree(C)
=
\Refine
\left(
C;
\Purpose,
\Authority,
\Evidence,
\Inv,
\Law,
\State,
\Verify,
\Prov
\right).
\]

The operator shall:

\begin{enumerate}
    \item preserve established canonical meaning;
    \item remove identified contradiction;
    \item remove unsupported assumptions where they are not
          authorized;
    \item resolve definitional conflation;
    \item expose unresolved obligations;
    \item preserve provenance;
    \item preserve authorized purpose;
    \item preserve invariant intent;
    \item preserve valid derivative relationships;
    \item introduce no unverified semantic claim as established.
\end{enumerate}

The operator therefore has two simultaneous constraints:

\[
\boxed{
\text{Defect Reduction}
}
\]

and:

\[
\boxed{
\text{Semantic Preservation}.
}
\]

A refinement that removes a defect by changing canonical meaning
without authorization is not a valid R³ refinement.

% ============================================================
\section{Fixed-Point Semantics}
% ============================================================

A canonical state $C^{*}$ is an \Rthree{} fixed point when:

\[
\boxed{
\Rthree(C^{*})=C^{*}
}
\]

with respect to the declared canonical comparison relation.

Because textual changes may occur without semantic change, the
fixed-point condition shall be evaluated semantically rather than
by raw character equality.

Define the canonical comparison relation:

\[
C_a\equiv_{\mathrm{can}} C_b
\]

when all materially relevant canonical properties are equivalent,
including:

\begin{enumerate}
    \item purpose;
    \item authority;
    \item definitions;
    \item invariants;
    \item law semantics;
    \item state semantics;
    \item verification semantics;
    \item provenance requirements;
    \item derivative constraints.
\end{enumerate}

Then the operational fixed-point condition is:

\[
\boxed{
\Rthree(C^{*})
\equiv_{\mathrm{can}}
C^{*}.
}
\]

This prevents stylistic or representational changes from
producing an artificial infinite refinement sequence.

% ============================================================
\section{Diminishing-Returns Boundary}
% ============================================================

Recursive refinement shall continue while a materially relevant
defect remains.

Let:

\[
\Delta(C_n,C_{n+1})
\]

denote the set of materially relevant canonical changes produced
by an R³ iteration.

The refinement process terminates when:

\[
\boxed{
\Delta(C_n,C_{n+1})=\varnothing
}
\]

with respect to the declared canonical comparison relation.

Further changes that affect only:

\begin{enumerate}
    \item typography;
    \item stylistic preference;
    \item non-semantic ordering;
    \item equivalent notation;
    \item non-material wording;
\end{enumerate}

do not constitute a reason to continue R³ refinement of the
canonical semantics.

Thus:

\[
\boxed{
\text{No Material Defect}
\Rightarrow
\text{Stop Refinement}.
}
\]

This defines the boundary between convergence and diminishing
returns.

% ============================================================
\section{Meta-Dynamical Mathematics}
% ============================================================

\MDM{} is the structural mathematical reference layer of the
canonical architecture.

It represents relationships among:

\[
\boxed{
\text{State},
\text{Transformation},
\text{Invariant},
\text{Observation},
\text{Mapping},
\text{Fixed Point}.
}
\]

A domain $D_i$ may be represented as:

\[
D_i=(S_i,O_i,F_i,I_i),
\]

where:

\begin{align*}
S_i &= \text{domain state space},\\
O_i &= \text{domain observations},\\
F_i &= \text{domain transformations},\\
I_i &= \text{domain invariants}.
\end{align*}

A common structural representation exists when a mapping:

\[
\phi_i:D_i\rightarrow M
\]

preserves the specified structures.

For a transformation:

\[
f_i:S_i\rightarrow S_i,
\]

a structure-preserving representation satisfies:

\[
\boxed{
\phi_i(f_i(x))
=
f_M(\phi_i(x))
}
\]

within the declared mathematical representation.

For an invariant:

\[
I_i(x)=\mathrm{true},
\]

the corresponding representation shall preserve the declared
invariant relation:

\[
\boxed{
I_i(x)
\Rightarrow
I_M(\phi_i(x)).
}
\]

The exact preservation conditions shall be specified for each
mapping.

MDM therefore provides the reference structure through which
different domains may be related without requiring their
substantive content to be identical.

% ============================================================
\section{Cross-Domain Structural Correspondence}
% ============================================================

\begin{definition}[Cross-Domain Isomorphism]
A cross-domain isomorphism is a structure-preserving mapping
between formally represented domains for which the specified
operations, relationships, and invariants are preserved under
the mapping.
\end{definition}

Cross-domain isomorphism does not imply:

\[
D_A=D_B.
\]

It instead establishes correspondence:

\[
\boxed{
D_A
\xleftrightarrow{\phi}
D_B
}
\]

with respect to the specified structure.

Therefore:

\[
\boxed{
\text{Different Domain Substance}
\neq
\text{Different Mathematical Structure}.
}
\]

The common MVL structural envelope is:

\[
\boxed{
\Purpose
\rightarrow
\Inv
\rightarrow
\State
\rightarrow
\Verify
\rightarrow
\delta.
}
\]

Domain-specific substance remains attached to the domain in which
it is authorized.

% ============================================================
\section{WAD-18 Deterministic Representation}
% ============================================================

\WAD{} is a deterministic fixed-point arithmetic representation
profile suitable for implementations requiring exact integer
semantics within its declared numerical domain.

A \WAD{} implementation may require:

\begin{enumerate}
    \item integer fixed-point representation;
    \item no floating-point arithmetic in normative or verdict
          modules;
    \item explicit overflow handling;
    \item explicit underflow handling where applicable;
    \item explicit division-by-zero handling;
    \item deterministic arithmetic;
    \item explicit state;
    \item fail-closed behavior for unresolved mandatory conditions.
\end{enumerate}

\WAD{} is not a necessary definition of \MVL{}.

The canonical dependency is:

\[
\boxed{
\Rthree{}
\rightarrow
\text{Invariant and Truth Pursuit}
\rightarrow
\text{Formal Representation}
\rightarrow
\text{Deterministic Implementation}.
}
\]

One possible implementation path is:

\[
\boxed{
\Rthree{}
\rightarrow
\MDM
\rightarrow
\WAD{}
\rightarrow
\text{Deterministic State Evaluation}.
}
\]

Other exact mathematical representations may conform when they
preserve the required canonical semantics.

% ============================================================
\section{Cognitive State Ledger}
% ============================================================

The Cognitive State Ledger (\CSL{}) provides a structured record
for state, evidence, reasoning status, verification, and
provenance.

A canonical ledger record may be represented as:

\[
\boxed{
\begin{gathered}
\CSL=
(
id,
purpose,
authority,
state,
evidence,
invariants,\\
operations,
result,
provenance,
version,
authorization,
obligations
)
\end{gathered}
}
\]

The ledger shall distinguish at minimum:

\begin{enumerate}
    \item established facts;
    \item derived propositions;
    \item assumptions;
    \item obligations;
    \item verification results;
    \item provenance records;
    \item superseded states.
\end{enumerate}

An obligation shall not be silently represented as an established
fact.

% ============================================================
\section{Human Intent and Custodial Authority}
% ============================================================

Machine verification does not independently determine the
normative purpose for which a law or system exists.

Human or otherwise legitimately authorized normative authority
therefore precedes machine formalization.

The canonical relationship is:

\[
\boxed{
\text{Authorized Human Purpose}
\rightarrow
\text{Intent Specification}
\rightarrow
\Rthree{}
\rightarrow
\Inv
\rightarrow
\Law
\rightarrow
\Verify.
}
\]

A Sovereign Custodian or Chief Custodial \& Invariant Officer
(\CCIO{}) may be defined as a custodial role responsible for
maintaining the authorized relationship among:

\[
\Purpose,\;
Authority,\;
\Inv,\;
\Prov,\;
\Law.
\]

The role may include:

\begin{enumerate}
    \item preserving canonical purpose;
    \item maintaining invariant definitions;
    \item maintaining authorization;
    \item resolving semantic ambiguity through authorized process;
    \item maintaining provenance;
    \item approving canonical revisions;
    \item ensuring derivative implementations remain semantically
          aligned.
\end{enumerate}

Custodial authority does not replace formal verification.

Formal verification does not replace legitimate normative
authority.

They are distinct layers of the architecture.

% ============================================================
\section{Machine-Verifiable Law Lifecycle}
% ============================================================

The canonical lifecycle is:

\[
\boxed{
\begin{array}{c}
\textbf{Purpose}\\
\downarrow\\
\textbf{Authority}\\
\downarrow\\
\textbf{Evidence}\\
\downarrow\\
\Rthree{}\\
\downarrow\\
\textbf{Invariant Extraction}\\
\downarrow\\
\MDM\\
\downarrow\\
\textbf{Formal Law}\\
\downarrow\\
\textbf{State Representation}\\
\downarrow\\
\textbf{Verification}\\
\downarrow\\
\textbf{Authorized Transition}\\
\downarrow\\
\textbf{Observation}\\
\downarrow\\
\Rthree{}\\
\downarrow\\
\textbf{Reverification}\\
\downarrow\\
\textbf{Truth-Constrained State}
\end{array}
}
\]

The lifecycle supports:

\begin{enumerate}
    \item creation;
    \item validation;
    \item formalization;
    \item testing;
    \item deployment;
    \item monitoring;
    \item revision;
    \item supersession;
    \item archival.
\end{enumerate}

No revision shall silently erase historical provenance.

% ============================================================
\section{Canonical Derivation Architecture}
% ============================================================

The canonical specification is the semantic source for all
authorized derivative artifacts.

\[
\boxed{
\Spec
\longrightarrow
\begin{cases}
\Spec_{\mathrm{Lean}}\\
\Spec_{\mathrm{JSONL}}\\
\Spec_{\mathrm{CodeOcean}}\\
\Spec_{\mathrm{Blockchain}}
\end{cases}
}
\]

Each derivative has a distinct representational function.

\begin{table}[h]
\centering
\begin{tabular}{p{0.22\textwidth}p{0.68\textwidth}}
\toprule
\textbf{Representation} & \textbf{Function}\\
\midrule
Lean 4 &
Formal representation of definitions, propositions, structures,
and proofs subject to the formalization.\\[0.5em]

JSONL &
Machine-readable serialization of canonical definitions,
relationships, examples, obligations, and metadata.\\[0.5em]

Code Ocean &
Executable reproducibility environment for defined computational
procedures, experiments, tests, and verification capsules.\\[0.5em]

Blockchain &
Cryptographic anchoring and provenance of canonical versions,
digests, attestations, and authorized state records.\\
\bottomrule
\end{tabular}
\caption{Authorized derivative representations}
\end{table}

No derivative becomes authoritative merely because it is
executable, immutable, computationally convenient, or externally
distributed.

Authority remains with the canonical semantic specification and
its authorized source.

% ============================================================
\section{Semantic Preservation}
% ============================================================

Let:

\[
\Sem(X)
\]

denote the semantic content of representation $X$ within its
declared scope.

For an exact derivative:

\[
\boxed{
\Sem(D_i)\equiv\Sem(\Spec)
}
\]

within the applicable representational domain.

If exact equivalence cannot be established, the relationship shall
be explicitly classified as one of:

\begin{enumerate}
    \item exact representation;
    \item structure-preserving representation;
    \item projection;
    \item approximation;
    \item unresolved mapping.
\end{enumerate}

An unresolved mapping shall not be represented as an established
equivalence.

% ============================================================
\section{Reproducibility}
% ============================================================

Let:

\[
J=(I,R,V,E)
\]

represent:

\begin{align*}
I &= \text{canonical inputs},\\
R &= \text{applicable rules},\\
V &= \text{versions},\\
E &= \text{declared execution conditions}.
\end{align*}

A reproducible verification process satisfies:

\[
\boxed{
J_1\equiv J_2
\Rightarrow
\Verify(J_1)=\Verify(J_2)
}
\]

within the declared equivalence conditions.

Reproducibility requires, where applicable:

\begin{enumerate}
    \item fixed input versions;
    \item fixed rule versions;
    \item explicit dependencies;
    \item explicit arithmetic semantics;
    \item explicit environmental assumptions;
    \item deterministic execution where required;
    \item preserved provenance;
    \item reproducible test procedures.
\end{enumerate}

A result that cannot be reproduced from its declared conditions
shall not be represented as fully reproducible.

% ============================================================
\section{Verification and Fail-Closed Semantics}
% ============================================================

For:

\[
\Inv=\{I_1,I_2,\ldots,I_n\},
\]

let:

\[
I_i(S)\in
\{
\mathrm{true},
\mathrm{false},
\mathrm{unknown}
\}.
\]

Then:

\[
\boxed{
\Verify(S,\Law)=\Accept
}
\]

only when every mandatory applicable invariant is established as
true and all other mandatory verification conditions are
satisfied.

If:

\[
\exists i:
I_i(S)=\mathrm{false},
\]

then:

\[
\boxed{
\Verify(S,\Law)=\Reject.
}
\]

If:

\[
\exists i:
I_i(S)=\mathrm{unknown},
\]

for a mandatory applicable condition, then:

\[
\boxed{
\Verify(S,\Law)=\Unknown.
}
\]

Thus:

\[
\boxed{
\Unknown\not\Rightarrow\Accept
}
\]

and:

\[
\boxed{
\Reject\not\equiv\Unknown.
}
\]

An affirmative violation and an inability to establish
compliance are distinct states.

% ============================================================
\section{Evidence and Provenance}
% ============================================================

The provenance chain may be represented as:

\[
\boxed{
\text{Source}
\rightarrow
\text{Authority}
\rightarrow
\text{Interpretation}
\rightarrow
\text{Formalization}
\rightarrow
\text{Implementation}
\rightarrow
\text{Execution}
\rightarrow
\text{Result}
}
\]

Each transformation shall be identifiable.

Where applicable, provenance shall contain:

\begin{enumerate}
    \item source identifier;
    \item authority;
    \item version;
    \item timestamp;
    \item transformation identifier;
    \item canonical digest;
    \item derivative digest;
    \item execution environment;
    \item verification result.
\end{enumerate}

A digest shall only be represented when actually computed.

No fabricated hash or provenance value is permissible.

% ============================================================
\section{Contradiction and Defect Handling}
% ============================================================

Contradictory conditions shall not be silently reconciled by
implementation preference.

If:

\[
I_a(S)=\mathrm{true}
\]

and:

\[
I_b(S)=\mathrm{false},
\]

then the state cannot simultaneously satisfy both invariants.

Where contradictory authoritative sources exist, the canonical
state shall preserve the contradiction until an authorized
resolution rule is established.

The canonical defect path is:

\[
\boxed{
\text{Defect}
\rightarrow
\text{Explicit State}
\rightarrow
\Rthree{}
\rightarrow
\text{Resolution}
\rightarrow
\text{Reverification}.
}
\]

This prevents implementation behavior from becoming an
undocumented substitute for semantic resolution.

% ============================================================
\section{Unknown and Obligation States}
% ============================================================

A canonical verification process shall distinguish:

\[
\boxed{
\text{TRUE},
\quad
\text{FALSE},
\quad
\text{UNKNOWN},
\quad
\Obligation
}
\]

where:

\[
\Obligation
\]

denotes a required condition that has not yet been established.

Therefore:

\[
\boxed{
\Obligation\neq\text{Truth}.
}
\]

Likewise:

\[
\boxed{
\Unknown\neq\text{Truth}.
}
\]

A non-consequential process may continue analysis while an
obligation remains unresolved.

A consequential acceptance shall require satisfaction of every
mandatory obligation applicable to that acceptance.

% ============================================================
\section{Formal Law and Machine Law}
% ============================================================

The transformation from normative source to machine-verifiable
predicate is:

\[
\boxed{
\begin{gathered}
\text{Normative Source}
\rightarrow
\text{Canonical Interpretation}
\rightarrow
\text{Formal Rule}
\rightarrow\\
\text{Invariant}
\rightarrow
\text{Executable Predicate}
\end{gathered}
}
\]

For a particular invariant:

\[
I:S\rightarrow\{\mathrm{true},\mathrm{false},\mathrm{unknown}\}.
\]

A deterministic implementation of $I$ may provide repeatable
evaluation.

However:

\[
\boxed{
\text{Deterministic Computation}
\neq
\text{Semantic Correctness}.
}
\]

Semantic correctness depends upon preservation of:

\[
\Purpose,\;
Authority,\;
\Evidence,\;
\Inv,\;
\Law,\;
\Prov.
\]

This distinction is foundational.

% ============================================================
\section{Legal Validity and Formal Verifiability}
% ============================================================

Machine-verifiable representation and legal validity are distinct
properties.

The architecture distinguishes:

\[
\boxed{
\text{Normative Authority}
}
\]

from:

\[
\boxed{
\text{Formal Representation}
}
\]

from:

\[
\boxed{
\text{Machine Verification}.
}
\]

Formalizing a rule does not independently create its legal
authority.

Likewise, a machine-verifiable predicate does not independently
establish that the underlying normative source was legitimately
authorized.

The canonical chain is therefore:

\[
\boxed{
\text{Authorized Normative Source}
\rightarrow
\text{Canonical Representation}
\rightarrow
\text{Formal Verification}
}
\]

rather than:

\[
\text{Machine Output}
\rightarrow
\text{Normative Authority}.
\]

This distinction preserves the boundary between law's source of
authority and the computational mechanism used to evaluate its
formal conditions.

% ============================================================
\section{Cross-Domain Law Envelope}
% ============================================================

Let:

\[
\Domain=\{D_1,D_2,\ldots,D_n\}
\]

be a collection of domains.

Each domain may define:

\[
\Law_i.
\]

The common MVL envelope is:

\[
\boxed{
\mathfrak{L}
=
\left\{
\Law_i
\mid
D_i\in\Domain
\right\}
}
\]

subject to common structural requirements:

\[
\Purpose,\;
Authority,\;
\Inv,\;
State,\;
\Verify,\;
Prov,\;
\delta.
\]

The envelope does not impose identical substantive rules.

It imposes a common representation and verification discipline.

Therefore:

\[
\boxed{
\text{Universal Structure}
+
\text{Domain-Specific Substance}
=
\text{Cross-Domain MVL}.
}
\]

% ============================================================
\section{Machine-Verifiable Safety}
% ============================================================

Safety may be represented as a formally defined invariant rather
than solely as an after-the-fact assessment.

For a safety invariant:

\[
I_S(S)=\mathrm{true}
\]

indicates that the specified safety condition is satisfied for
state $S$.

A Safety Asset Class (\SAC{}) representation may classify
verified safety properties as explicit attributes of a state.

The canonical relationship is:

\[
\boxed{
\Purpose
\rightarrow
I_S
\rightarrow
\Verify
\rightarrow
\text{Verified Safety State}.
}
\]

Safety verification remains distinct from economic value,
normative authority, or other classifications unless an explicit
canonical relationship is established.

% ============================================================
\section{Security and Adversarial Integrity}
% ============================================================

An MVL implementation shall account for integrity conditions
including, where applicable:

\begin{enumerate}
    \item hidden state;
    \item altered provenance;
    \item version substitution;
    \item numerical overflow;
    \item numerical underflow;
    \item division by zero;
    \item malformed input;
    \item ambiguous interpretation;
    \item unauthorized rule modification;
    \item stale rule versions;
    \item dependency substitution;
    \item environment drift;
    \item inconsistent state synchronization.
\end{enumerate}

Where a mandatory integrity condition cannot be established:

\[
\boxed{
\text{Integrity Unknown}
\Rightarrow
\Verify\neq\Accept.
}
\]

% ============================================================
\section{Canonical Versioning and Supersession}
% ============================================================

Every canonical release shall have an immutable version
identifier.

A canonical release may be represented as:

\[
C_v=(id,v,h,t,a),
\]

where:

\begin{align*}
id &= \text{canonical identity},\\
v &= \text{version},\\
h &= \text{canonical digest},\\
t &= \text{publication or activation time},\\
a &= \text{authorized authority}.
\end{align*}

A subsequent version:

\[
C_{v+1}
\]

shall preserve a traceable relationship to:

\[
C_v.
\]

Supersession shall not erase historical provenance.

A new version shall identify whether its change is:

\begin{enumerate}
    \item semantic;
    \item structural;
    \item representational;
    \item editorial;
    \item corrective.
\end{enumerate}

% ============================================================
\section{Blockchain Provenance}
% ============================================================

Blockchain technology may provide an external cryptographic
provenance layer.

The relationship is:

\[
\boxed{
\text{Canonical State}
\rightarrow
\text{Canonical Digest}
\rightarrow
\text{Blockchain Record}.
}
\]

The blockchain record may establish evidence that a specified
digest was anchored in a specified ledger state.

It does not independently establish semantic correctness.

Therefore:

\[
\boxed{
\text{Blockchain Immutability}
\neq
\text{Semantic Truth}.
}
\]

Blockchain anchoring is consequently a provenance function, not
the canonical source of normative meaning.

% ============================================================
\section{Lean 4 Formalization}
% ============================================================

Lean 4 may serve as a formal representation environment for
canonical definitions, structures, propositions, and proofs.

The derivation is:

\[
\boxed{
\Spec
\rightarrow
\Spec_{\mathrm{Lean}}.
}
\]

The canonical relationship is:

\[
\text{Canonical Definition}
\rightarrow
\text{Formal Proposition}
\rightarrow
\text{Formal Proof}.
\]

Where a canonical proposition cannot be formalized or proven under
the declared assumptions, the discrepancy shall be represented
explicitly.

A failed proof shall not be represented as a successful proof.

A proof of a formalized proposition shall not automatically be
interpreted as proof of an unstated proposition outside that
formalization.

% ============================================================
\section{JSONL Canonical Representation}
% ============================================================

JSONL may provide a machine-readable serialization of canonical
semantic objects.

A canonical JSONL record should preserve, where applicable:

\begin{enumerate}
    \item canonical identifier;
    \item object type;
    \item canonical definition;
    \item dependencies;
    \item invariants;
    \item authority;
    \item provenance;
    \item version;
    \item verification status;
    \item obligations;
    \item relationships to other canonical records.
\end{enumerate}

The JSONL representation shall be derived from the canonical
semantic root.

It shall not introduce a new normative definition merely because
the serialization format makes such a definition convenient.

% ============================================================
\section{Code Ocean Reproducibility Capsules}
% ============================================================

A Code Ocean capsule or equivalent executable environment may
serve as a reproducibility realization of a defined canonical
computational procedure.

The capsule shall identify, where applicable:

\begin{enumerate}
    \item source code;
    \item dependencies;
    \item runtime;
    \item inputs;
    \item tests;
    \item expected results;
    \item actual results;
    \item version;
    \item provenance.
\end{enumerate}

The capsule is therefore:

\[
\boxed{
\text{Executable Evidence of a Defined Canonical Procedure}.
}
\]

An executable result shall not be generalized beyond its declared
test domain without additional evidence.

% ============================================================
\section{Canonical Conformance}
% ============================================================

An implementation conforms to this specification within a
declared scope when it demonstrates:

\begin{enumerate}
    \item explicit purpose;
    \item explicit authority;
    \item explicit invariants;
    \item defined state;
    \item defined verification semantics;
    \item provenance;
    \item reproducibility;
    \item semantic preservation;
    \item fail-closed handling of unresolved mandatory conditions;
    \item version control;
    \item auditable derivation.
\end{enumerate}

A conformance claim shall identify:

\begin{enumerate}
    \item the canonical version;
    \item the derivative representation;
    \item the applicable domain;
    \item the verification scope;
    \item the tested conditions;
    \item any unresolved obligations.
\end{enumerate}

No implementation shall claim universal conformance solely from
conformance within a limited computational domain.

% ============================================================
\section{Canonical Derivation Rule}
% ============================================================

The central derivation rule is:

\[
\boxed{
\parbox{0.85\textwidth}{\centering\textbf{
Every derivative representation shall be derivable from a
canonical semantic specification whose purpose, authority,
invariants, provenance, state model, and verification semantics
are explicit.
}}
}
\]

Accordingly:

\[
\boxed{
\Spec
\rightarrow
\Spec_{\mathrm{Lean}}
}
\]

\[
\boxed{
\Spec
\rightarrow
\Spec_{\mathrm{JSONL}}
}
\]

\[
\boxed{
\Spec
\rightarrow
\Spec_{\mathrm{CodeOcean}}
}
\]

\[
\boxed{
\Spec
\rightarrow
\Spec_{\mathrm{Blockchain}}
}
\]

These are controlled derivations rather than independent sources
of meaning.

% ============================================================
\section{Derivative Reconciliation Loop}
% ============================================================

Every derivative representation may be evaluated against the
canonical semantic root.

The reconciliation loop is:

\[
\boxed{
\Spec
\rightarrow
D_i
\rightarrow
\text{Derivative Verification}
\rightarrow
\begin{cases}
\text{Conformant}\\
\text{Defective}\\
\text{Unresolved}
\end{cases}
}
\]

If a derivative is defective or unresolved:

\[
\boxed{
D_i
\rightarrow
\Rthree
\rightarrow
\Spec'
\rightarrow
D_i'
}
\]

shall be used to reconcile the discrepancy.

A derivative shall not independently rewrite the canonical source
in order to conceal a discrepancy.

If the discrepancy reveals an error in the canonical semantic
root, the canonical root may be revised through the authorized
R³ process.

This creates the controlled refinement relation:

\[
\boxed{
\Spec
\leftrightarrow
\text{Formal}
\leftrightarrow
\text{Executable}
\leftrightarrow
\text{Observed}
}
\]

with \Rthree{} governing semantic reconciliation.

% ============================================================
\section{Canonical Verification Loop}
% ============================================================

The complete architecture is:

\[
\boxed{
\begin{aligned}
&\textbf{BEGIN WITH PURPOSE}\\
&\downarrow\\
&\text{Establish Authority}\\
&\downarrow\\
&\text{Collect Evidence}\\
&\downarrow\\
&\Rthree{}\\
&\downarrow\\
&\text{Identify Invariants}\\
&\downarrow\\
&\MDM\\
&\downarrow\\
&\text{Formalize Law}\\
&\downarrow\\
&\text{Represent State}\\
&\downarrow\\
&\text{Verify}\\
&\downarrow\\
&\text{Authorize Transition}\\
&\downarrow\\
&\text{Observe}\\
&\downarrow\\
&\Rthree{}\\
&\downarrow\\
&\text{Resolve Material Defects}\\
&\downarrow\\
&\text{Reverify}\\
&\downarrow\\
&\boxed{\textbf{END IN TRUTH}}
\end{aligned}
}
\]

The architecture is recursive because observation can expose
defects that were not established at an earlier state.

The recursive process terminates at the declared fixed-point
condition:

\[
\boxed{
\Rthree(C^{*})
\equiv_{\mathrm{can}}
C^{*}
}
\]

subject to the declared verification domain.

% ============================================================
\section{Canonical Mathematical Closure}
% ============================================================

The complete canonical object may be represented as:

\[
\boxed{
C=
(
\Purpose,
\Authority,
\Evidence,
\Inv,
\Law,
\State,
\Verify,
\Prov,
v
)
}
\]

and the refinement operator as:

\[
\boxed{
\Rthree:
\mathfrak{C}\rightarrow\mathfrak{C}.
}
\]

The canonical fixed point is:

\[
\boxed{
C^{*}
=
\Rthree(C^{*})
}
\]

under canonical semantic equivalence.

The truth-constrained state produced by the specification is:

\[
\boxed{
S^{*}
\in
\State
}
\]

such that:

\[
\boxed{
\Verify(S^{*},\Law)=\Accept
}
\]

within the declared verification domain.

The two fixed points shall not be conflated:

\[
\boxed{
\text{Canonical Fixed Point}
\neq
\text{Verified Runtime State}.
}
\]

The former concerns semantic specification convergence.

The latter concerns admissibility of a particular state.

% ============================================================
\section{Foundational Principles}
% ============================================================

The specification is governed by the following principles:

\begin{enumerate}
    \item \textbf{Purpose precedes consequential computation.}
    \item \textbf{Authorized purpose precedes normative formalization.}
    \item \textbf{\Rthree{} recursively pursues absolute truth.}
    \item \textbf{\MDM{} provides the structural mathematical reference.}
    \item \textbf{Explicit invariants replace undocumented assumptions.}
    \item \textbf{Provenance preserves the origin and continuity of meaning.}
    \item \textbf{Reproducibility makes verification independently examinable.}
    \item \textbf{Determinism applies where deterministic verification is specified.}
    \item \textbf{Unknown is not acceptance.}
    \item \textbf{Rejection and non-verification are distinct states.}
    \item \textbf{No derivative may silently alter canonical meaning.}
    \item \textbf{No implementation may silently substitute its own authority.}
    \item \textbf{Cross-domain structure does not require identical substantive rules.}
    \item \textbf{Machine verification does not independently create normative authority.}
    \item \textbf{Canonical convergence is semantic, not merely textual.}
    \item \textbf{Refinement stops when no material canonical defect remains.}
\end{enumerate}

% ============================================================
\section{Canonical Integrity Rules}
% ============================================================

The following rules constitute the minimum semantic integrity
conditions of the architecture.

\begin{enumerate}
    \item
    \[
    \boxed{
    \text{No derivative may contain more authority than its source.}
    }
    \]

    \item
    \[
    \boxed{
    \text{No implementation may silently change canonical meaning.}
    }
    \]

    \item
    \[
    \boxed{
    \text{No unverified condition may be represented as verified.}
    }
    \]

    \item
    \[
    \boxed{
    \text{No unresolved contradiction may be silently converted into a rule.}
    }
    \]

    \item
    \[
    \boxed{
    \text{No fabricated provenance may be represented as evidence.}
    }
    \]

    \item
    \[
    \boxed{
    \text{No semantic convergence claim may be made without a declared comparison relation.}
    }
    \]
\end{enumerate}

% ============================================================
\section{Final Canonical Statement}
% ============================================================

Machine-Verifiable Law is the formal representation of
authorized normative intent, applicable law, invariants,
evidence, provenance, states, and authorized transitions in a
form that permits explicit, reproducible, auditable, and
appropriately deterministic verification within a declared
domain.

Its governing architecture is:

\[
\boxed{
\Purpose
\rightarrow
\Authority
\rightarrow
\Rthree{}
\rightarrow
\MDM
\rightarrow
\Inv
\rightarrow
\Law
\rightarrow
\State
\rightarrow
\Verify
\rightarrow
\Truth
}
\]

Its governing semantic-root principle is:

\[
\boxed{
\begin{gathered}
\textbf{ONE CANONICAL SEMANTIC ROOT}\\
\downarrow\\
\textbf{MANY CONTROLLED REPRESENTATIONS}
\end{gathered}
}
\]

Its governing refinement principle is:

\[
\boxed{
\textbf{
R}^{3}
=
\textbf{RECURSIVE PURSUIT OF ABSOLUTE TRUTH}
}
\]

Its governing mathematical principle is:

\[
\boxed{
\textbf{
MDM PROVIDES THE COMMON STRUCTURAL REFERENCE
}
}
\]

Its governing integrity principle is:

\[
\boxed{
\textbf{
NO SILENT SEMANTIC DRIFT
}
}
\]

Its governing verification principle is:

\[
\boxed{
\textbf{
NO VERIFIED STATE WITHOUT VERIFIABLE CONDITIONS
}
}
\]

Its governing convergence condition is:

\[
\boxed{
\boldsymbol{
\Rthree(C^{*})
\equiv_{\mathrm{can}}
C^{*}
}
}
\]

and its governing direction is:

\[
\boxed{
\text{\Large\textbf{BEGIN WITH PURPOSE. END IN TRUTH.}}
}
\]

\vfill

\begin{center}
\textit{
Purpose establishes direction.\\
Authority establishes normative source.\\
R\textsuperscript{3} establishes recursive pursuit.\\
MDM establishes structural mathematics.\\
Invariants establish admissibility.\\
Formalization establishes representation.\\
Verification establishes state.\\
Provenance establishes continuity.\\
Reproducibility establishes independent examinability.\\
Truth establishes the fixed point.
}
\end{center}

\end{document}
